特征值分解很有力量，却只在方阵上直接成立，而且只有在矩阵拥有足够多特征向量时才真正好用。现实中的数据矩阵往往是长方形的：每一行是一个样本，每一列是一项特征。即使矩阵恰好是方阵，它也未必对称、未必能被一组正交特征向量对角化。

奇异值分解解决的正是这个缺口。它告诉我们，任何线性映射都可以拆成"选择正交的输入方向---分别伸缩---转到正交的输出方向"。这一个结构同时解释了伪逆、最小二乘的敏感性、低秩压缩和主成分分析。

\subsection{先看一个例子}

取 \(A=\begin{bmatrix}3&0\\4&5\end{bmatrix}\)。试一下 \(A^{\trans}A=\begin{bmatrix}25&20\\20&25\end{bmatrix}\)，特征值 \(45\) 和 \(5\)，所以奇异值 \(\sigma_1=\sqrt{45}=3\sqrt5\approx6.708\)，\(\sigma_2=\sqrt5\approx2.236\)。对应右奇异向量 \(v_1=(1,1)^{\trans}/\sqrt2\)，\(v_2=(1,-1)^{\trans}/\sqrt2\)，左奇异向量 \(u_1=Av_1/\sigma_1\)，\(u_2=Av_2/\sigma_2\)。单位圆在 \(A\) 作用后变成一个椭圆，长半轴长度 \(\sigma_1\) 沿 \(u_1\) 方向，短半轴长度 \(\sigma_2\) 沿 \(u_2\) 方向。两个奇异值都不为零，所以 \(A\) 可逆；若有一个零奇异值，就说明存在某些输入方向被完全压成零，矩阵非满秩。

建议按下面的顺序连续阅读：

\begin{enumerate}
\tightlist
\item
  任意矩阵的奇异值分解：从单位球的像出发，推导 \(A=U\Sigma V^{\trans}\)。
\item
  伪逆与病态问题：哪些信息可以反解，哪些微小误差会被放大。
\item
  低秩近似：为什么截断 SVD 是最好的秩约束近似。
\item
  从数据矩阵到 PCA：把前面的结构翻译成降维与数据解释。
\end{enumerate}

前置内容是对称矩阵的谱定理和最小二乘。读完本单元后，Fourier 分析会显得很自然：它也是寻找一组能把结构化线性变换变简单的正交坐标。

\subsection{怎么读这一章}

先把 SVD 看成"旋转---按正交方向缩放---再旋转"，随后处理不可逆或病态系统的伪逆，再用截断奇异值获得最佳低秩近似，最后把中心化数据矩阵的 SVD 与 PCA 连接起来。

\subsection{本章篇目}

以下篇目按概念依赖排列；若从具体问题进入，也可以先打开最接近当前困惑的一篇，再沿篇内链接回补。

\begin{itemize}
\tightlist
\item \hyperref[sec:03-Linear-Algebra/07-SVD-and-Data-Applications/01]{``任意矩阵的奇异值分解''}；
\item \hyperref[sec:03-Linear-Algebra/07-SVD-and-Data-Applications/02]{``伪逆与病态问题''}；
\item \hyperref[sec:03-Linear-Algebra/07-SVD-and-Data-Applications/03]{``低秩近似：用少数方向保留主要结构''}；
\item \hyperref[sec:03-Linear-Algebra/07-SVD-and-Data-Applications/04]{``从数据矩阵到 PCA''}。
\end{itemize}
